% --- Part 1: Introduction & Motivation (Slides 1-3) ---
\section{Introducción y Motivación}

\begin{frame}{Título y Autores}
    \begin{center}
        \Huge \textbf{Librería Python para CKKS en GPU} \\
        \large \textit{Operaciones Tensoriales con Cifrado Homomórfico} \\
        \vspace{1cm}
        \normalsize \textbf{Rafael Eduardo Contreras} \\
        \textit{Universidad Central de Venezuela \\ Facultad de Ciencias}
    \end{center}
\end{frame}

\begin{frame}{El Dilema de la Privacidad}
    \begin{columns}
        \column{0.6\textwidth}
        \begin{itemize}
            \item Explosión del \textbf{Machine Learning (ML)} en áreas sensibles.
            \item \textbf{Sectores:} Medicina, Finanzas, etc.
            \item \textbf{Riesgo:} Fuga de datos al usar servicios de terceros en la nube para inferencia.
        \end{itemize}
        
        \column{0.4\textwidth}
        \begin{figure}
            \centering
            \includegraphics[width=\textwidth]{fhe_on_ml.png}
        \end{figure}
    \end{columns}
\end{frame}

\begin{frame}{Cifrado Homomórfico (HE) como Solución}
    \begin{itemize}
        \item \textbf{Definición:} Computar directamente sobre datos cifrados.
        \item \textbf{Propiedad:} El resultado, una vez descifrado, coincide con el obtenido en texto plano.
        \item \textbf{Beneficio:} El servidor nunca tiene acceso a la información sensible.
    \end{itemize}
\end{frame}
